\documentclass[12pt]{article}

\usepackage{amsmath}
\usepackage{amssymb}
\usepackage{graphicx}
\usepackage{hyperref}
\usepackage[margin=1in]{geometry}
\usepackage{natbib}
\usepackage{xcolor}
\usepackage{tcolorbox}
\usepackage{booktabs}
\usepackage{tabularx}
\usepackage{lineno}
\linenumbers
\tcbuselibrary{breakable,skins}
\bibliographystyle{aasjournal}

\newcommand{\bk}{\mathbf{k}}
\newcommand{\bx}{\mathbf{x}}
\newcommand{\eV}{\,\mathrm{eV}}
\newcommand{\kpc}{\,\mathrm{kpc}}
\newcommand{\kms}{\,\mathrm{km\,s^{-1}}}
\newcommand{\Lag}{\mathcal{L}}
\newcommand{\Ham}{\mathcal{H}}

\title{Field Theory Foundations for Scalar Dark Matter:\\
       From the Principle of Least Action to the Classical Wave Limit}
\author{}
\date{\today}

\begin{document}

\maketitle
\tableofcontents
\newpage

\begin{tcolorbox}[colback=yellow!10!white, colframe=orange!75!black,
                  title={\bfseries Disclaimer}]
These notes were compiled by an AI assistant (Claude, Anthropic) and have
\textbf{not been reviewed by a human domain expert}.  They may contain
errors in derivations, physical reasoning, or numerical estimates.
Cross-check key equations against the primary literature before use.
\end{tcolorbox}
\vspace{0.5em}

%----------------------------------------------------------------------
\section{Purpose and scope}
\label{sec:intro}

The fuzzy-dark-matter notes (\texttt{fdm.tex}) start from the action of
a real scalar field minimally coupled to gravity,
%
\begin{equation}
  S = \int d^4x\,\sqrt{-g}\left[
       -\tfrac{1}{2}\,g^{\mu\nu}\,\partial_\mu\phi\,\partial_\nu\phi
       -\tfrac{1}{2}\,m^2\phi^2
      \right],
  \label{eq:scalar_action}
\end{equation}
%
and derive everything else --- misalignment, Schr\"odinger--Poisson,
solitons --- from it.  A short appendix there unpacks the symbols.
These notes expand that appendix into a self-contained account of the
field-theory background: where actions like \eqref{eq:scalar_action}
come from, what happens when the field is quantised, and --- the payoff
for dark-matter physics --- why a quantum field with occupation numbers
of order $10^{96}$ per de Broglie volume may be treated as a
\emph{classical} wave.

Nothing here is FDM-specific.  The material is the standard scalar
field theory found in the first chapters of any QFT course
\citep{Tong2007, Coleman2011, Srednicki2007, PeskinSchroeder1995};
these notes select the subset needed to read the cosmology literature
on ultralight scalars \citep{Hui2017, Marsh2016} and present it for
readers whose background is astrophysics or numerical cosmology rather
than particle theory.

Roadmap: \S\ref{sec:units} fixes units.
\S\ref{sec:action}--\ref{sec:free} develop \emph{classical} field
theory: the action principle, relativistic notation, and the free
scalar field.  \S\ref{sec:noether}--\ref{sec:hamiltonian} cover
symmetries, conserved currents, and the Hamiltonian formulation ---
the bridge to quantisation.  \S\ref{sec:quantization} quantises the
field; \S\ref{sec:classical} explains the classical-wave limit via
coherent states.  \S\ref{sec:interactions} adds interactions and the
axion potential, \S\ref{sec:curved} puts the field on curved
spacetime, and \S\ref{sec:conventions} collects the convention traps.

%----------------------------------------------------------------------
\section{Natural units}
\label{sec:units}

Particle physics sets $\hbar = c = 1$.  Then energy, mass, momentum,
inverse length, and inverse time all carry the same unit,
conventionally electron-volts:
%
\begin{equation}
  [E] = [m] = [p] = [x^{-1}] = [t^{-1}] = \mathrm{eV}.
\end{equation}
%
Conversions use $\hbar c = 197.3\;\mathrm{eV\,nm} =
1.973\times10^{-7}\;\mathrm{eV\,m}$ and $\hbar =
6.582\times10^{-16}\;\mathrm{eV\,s}$.

Two FDM-flavoured examples:
%
\begin{itemize}
  \item \textbf{Compton wavelength.}  For $m = 10^{-22}\eV$,
        $\lambda_C \equiv \hbar/(mc) = \hbar c/(mc^2)
        = 1.97\times10^{15}\,\mathrm{m} \simeq 0.064\,\mathrm{pc}$.
        This is the scale below which relativistic field dynamics
        matter; FDM structure formation happens far above it.
  \item \textbf{de Broglie wavelength.}  For a virial velocity
        $v = 10\kms$,
        $\lambda_{\rm dB}/2\pi = \hbar/(mv) = \lambda_C\,(c/v)
        \simeq 1.9\kpc$ --- the famous kpc-scale wave length that
        motivates the whole subject \citep{Hui2017}.
\end{itemize}
%
Factors of $\hbar$ and $c$ are restored at the end of a calculation by
dimensional analysis: insert the unique combination of $\hbar$ and $c$
that gives the target unit.

%----------------------------------------------------------------------
\section{The principle of least action}
\label{sec:action}

\subsection{Mechanics}
\label{sec:action_mech}

Given a Lagrangian $L(q,\dot q,t)$ for a system with one generalised
coordinate $q(t)$, the action is
%
\begin{equation}
  S[q] = \int_{t_1}^{t_2}\! L(q(t),\dot q(t), t)\,dt.
  \label{eq:action_mech}
\end{equation}
%
The physical trajectory makes $S$ stationary under variations
$q(t)\to q(t) + \delta q(t)$ with fixed endpoints
$\delta q(t_1) = \delta q(t_2) = 0$.  Compute $\delta S$ by chain
rule:
%
\begin{equation}
  \delta S
    = \int_{t_1}^{t_2}\!
        \left[
          \frac{\partial L}{\partial q}\,\delta q
          + \frac{\partial L}{\partial\dot q}\,\delta\dot q
        \right] dt.
\end{equation}
%
Use $\delta\dot q = \frac{d}{dt}\delta q$ (variation and time
derivative commute) and integrate the second term by parts:
%
\begin{equation}
  \int_{t_1}^{t_2}\!\frac{\partial L}{\partial\dot q}\,
    \frac{d(\delta q)}{dt}\,dt
   = \underbrace{\left[\frac{\partial L}{\partial\dot q}\,\delta q\right]_{t_1}^{t_2}}_{=0\;\text{by BCs}}
     - \int_{t_1}^{t_2}\!\frac{d}{dt}\!\left(\frac{\partial L}{\partial\dot q}\right)\,\delta q\,dt.
\end{equation}
%
Requiring $\delta S = 0$ for arbitrary $\delta q(t)$ forces the
integrand to vanish pointwise (variational lemma), giving the
\textbf{Euler--Lagrange equation}
%
\begin{equation}
  \frac{d}{dt}\!\frac{\partial L}{\partial\dot q}
    - \frac{\partial L}{\partial q} = 0.
  \label{eq:EL_mech}
\end{equation}
%
Sanity check: $L = \tfrac{1}{2}m\dot q^2 - V(q)$ gives
$m\ddot q = -V'(q)$, Newton's second law.

Why formulate mechanics this way at all?  Three reasons that carry
over to fields: (i) the action is a single scalar, so symmetries are
manifest --- if $S$ is invariant under a transformation, the equations
of motion are too; (ii) constraints and coordinate changes are easy
(no force bookkeeping); (iii) the action is the quantity that appears
in the path integral, so the classical$\to$quantum step is shortest
from here.

\subsection{Fields}
\label{sec:action_field}

Replace the finite set of coordinates $\{q_i(t)\}$ by a field
$\phi(\bx, t)$: one degree of freedom \emph{per spatial point}.  The
Lagrangian becomes a spatial integral of a local \emph{Lagrangian
density} $\Lag$:
%
\begin{equation}
  L(t) = \int d^3\bx\,\Lag\bigl(\phi,\,\partial_\mu\phi\bigr),
  \qquad
  S[\phi] = \int dt\,L(t) = \int d^4x\,\Lag.
  \label{eq:action_field}
\end{equation}
%
Locality --- $\Lag$ depends on $\phi$ and finitely many derivatives at
a single point --- is a physical postulate: it encodes the absence of
action-at-a-distance and, together with Lorentz invariance, leads to
causal propagation.

Varying $\phi(x)\to\phi(x)+\delta\phi(x)$ with $\delta\phi$ vanishing
on the spacetime boundary and integrating by parts with the 4D
divergence theorem gives, in exact analogy to the mechanics case, the
\textbf{field-theory Euler--Lagrange equation}
%
\begin{equation}
  \partial_\mu\!\frac{\partial\Lag}{\partial(\partial_\mu\phi)}
    - \frac{\partial\Lag}{\partial\phi} = 0.
  \label{eq:field_EL}
\end{equation}
%
The mechanics time derivative $\frac{d}{dt}$ has become a spacetime
divergence $\partial_\mu$; everything else is unchanged.

%----------------------------------------------------------------------
\section{Relativistic notation}
\label{sec:relativity}

\subsection{Index notation and Einstein summation}
\label{sec:index}

Spacetime points have four coordinates
$x^\mu = (t, \bx) = (x^0, x^1, x^2, x^3)$, $\mu\in\{0,1,2,3\}$.
Repeated upper--lower index pairs are summed (\emph{Einstein
convention}):
%
\begin{equation}
  A^\mu B_\mu \equiv \sum_{\mu=0}^{3} A^\mu B_\mu.
\end{equation}
%
\textbf{Contravariant} components $A^\mu$ (index up) belong to
vectors ($x^\mu$, $p^\mu$); \textbf{covariant} components $A_\mu$
(index down) belong to covectors (1-forms), the prototypical example
being the gradient
%
\begin{equation}
  \partial_\mu \equiv \frac{\partial}{\partial x^\mu}
    = \left(\frac{\partial}{\partial t},\,\nabla\right).
\end{equation}
%
The metric raises and lowers indices, $A_\mu = g_{\mu\nu}A^\nu$,
$A^\mu = g^{\mu\nu}A_\nu$, with $g^{\mu\nu}$ the matrix inverse of
$g_{\mu\nu}$.  Fully contracted expressions such as
$g^{\mu\nu}\partial_\mu\phi\,\partial_\nu\phi$ are scalars ---
invariant under coordinate changes.  This is the whole point of the
notation: any expression with all indices contracted in upper--lower
pairs is automatically frame-independent.

\subsection{The metric and $\sqrt{-g}$}
\label{sec:metric}

The metric tensor $g_{\mu\nu}(x)$ defines the invariant interval
between neighbouring events,
%
\begin{equation}
  ds^2 = g_{\mu\nu}\,dx^\mu\,dx^\nu.
\end{equation}
%
Flat (Minkowski) spacetime in the ``mostly plus'' signature used
throughout these notes:
%
\begin{equation}
  g_{\mu\nu} = \eta_{\mu\nu} = \mathrm{diag}(-1, +1, +1, +1),
  \qquad
  ds^2 = -dt^2 + d\bx^2.
\end{equation}
%
The flat FRW background of cosmology has
$g_{\mu\nu} = \mathrm{diag}(-1, a^2, a^2, a^2)$ with scale factor
$a(t)$; comoving spatial derivatives pick up $a^{-1}$ factors when
converted to physical ones.

The quantity $\sqrt{-g} \equiv \sqrt{-\det(g_{\mu\nu})}$ makes
$d^4x\,\sqrt{-g}$ a coordinate-invariant volume element: under a
coordinate change the Jacobian of $d^4x$ is cancelled by the
transformation of $\sqrt{-g}$.  In Minkowski $\sqrt{-g}=1$; in FRW
$\sqrt{-g}=a^3$, which is where the Hubble friction term in the
cosmological Klein--Gordon equation comes from
(\S\ref{sec:curved}).

\subsection{Scalar fields and Lorentz invariance}
\label{sec:scalar}

A \emph{scalar field} $\phi(x)$ assigns one real number to every
spacetime point and transforms trivially under Lorentz
transformations: $\phi'(x') = \phi(x)$.  It is the simplest
relativistic field (contrast vectors $A^\mu$, spinors $\psi_\alpha$,
tensors $h_{\mu\nu}$).  Its gradient $\partial_\mu\phi$ is a
covector, with Lorentz-invariant square
%
\begin{equation}
  g^{\mu\nu}\,\partial_\mu\phi\,\partial_\nu\phi
    = -\dot\phi^2 + (\nabla\phi)^2
  \qquad \text{(Minkowski)}.
  \label{eq:kinetic_square}
\end{equation}
%
This combination is the relativistic generalisation of $\dot q^2$:
time and space derivatives enter together, with opposite signs fixed
by the signature.

%----------------------------------------------------------------------
\section{The free scalar field}
\label{sec:free}

\subsection{Anatomy of the action}
\label{sec:anatomy}

The flat-space free scalar action, term by term:
%
\begin{equation*}
  S = \underbrace{\int d^4x}_{\text{spacetime integral}}
  \Bigl[
  \underbrace{-\tfrac{1}{2}\,\eta^{\mu\nu}\partial_\mu\phi\,\partial_\nu\phi}_{\text{kinetic term}}
  \;
  \underbrace{-\;\tfrac{1}{2}m^2\phi^2}_{\text{mass term}}
  \Bigr]
  = \int d^4x \left[ \tfrac{1}{2}\dot\phi^2 - \tfrac{1}{2}(\nabla\phi)^2
      - \tfrac{1}{2}m^2\phi^2 \right].
\end{equation*}
%
With the signature contract of Eq.~\eqref{eq:kinetic_square}, the
overall minus sign on the kinetic term gives the time derivative a
\emph{positive} coefficient --- the canonical kinetic energy
$\tfrac{1}{2}\dot\phi^2$ --- while spatial gradients and the mass term
enter as potential-energy costs.  The Lagrangian density is exactly
``$T - V$'' per unit volume, with three contributions:
%
\begin{center}
\begin{tabular}{@{}lll@{}}
\toprule
Term & Mechanics analogue & Penalises \\
\midrule
$\tfrac{1}{2}\dot\phi^2$        & kinetic energy      & --- \\
$\tfrac{1}{2}(\nabla\phi)^2$    & spring coupling to neighbours & spatial variation \\
$\tfrac{1}{2}m^2\phi^2$         & harmonic restoring force      & displacement from $\phi=0$ \\
\bottomrule
\end{tabular}
\end{center}
%
The picture to keep: a scalar field is an infinite lattice of coupled
harmonic oscillators, one at each point of space, in the continuum
limit.  Everything that follows --- plane-wave modes, quantisation,
occupation numbers --- is the direct field-theory transcription of
coupled-oscillator physics.

\subsection{Klein--Gordon equation and plane waves}
\label{sec:kg}

Apply Eq.~\eqref{eq:field_EL}:
%
\begin{equation}
  \frac{\partial\Lag}{\partial\phi} = -m^2\phi,
  \qquad
  \frac{\partial\Lag}{\partial(\partial_\mu\phi)}
    = -\eta^{\mu\nu}\partial_\nu\phi = -\partial^\mu\phi,
\end{equation}
%
so
%
\begin{equation}
  \Box\phi = m^2\phi,
  \qquad
  \Box \equiv \partial_\mu\partial^\mu = -\partial_t^2 + \nabla^2,
  \label{eq:KG_flat}
\end{equation}
%
the \textbf{Klein--Gordon (KG) equation}.  Insert a plane wave
$\phi \propto e^{i(\bk\cdot\bx - \omega t)}$:
%
\begin{equation}
  \omega^2 = k^2 + m^2.
  \label{eq:dispersion}
\end{equation}
%
This is just the relativistic energy--momentum relation
$E^2 = p^2 + m^2$ with $E=\hbar\omega$, $p=\hbar k$ --- the first
hint that the field's excitations are particles of mass $m$.

Two limits of \eqref{eq:dispersion}:
%
\begin{itemize}
  \item \textbf{Relativistic} ($k \gg m$): $\omega \simeq k$, waves
        propagate at the speed of light.
  \item \textbf{Non-relativistic} ($k \ll m$):
        $\omega \simeq m + k^2/2m$.  After factoring out the fast
        phase $e^{-imt}$, the residual envelope obeys a
        Schr\"odinger equation --- this is the reduction performed in
        \texttt{fdm.tex} \S3, and the group velocity
        $v_g = d\omega/dk = k/m$ is the particle velocity.
\end{itemize}

%----------------------------------------------------------------------
\section{Symmetries and Noether's theorem}
\label{sec:noether}

\subsection{The theorem}
\label{sec:noether_thm}

Suppose an infinitesimal transformation
$\phi \to \phi + \epsilon\,\Delta\phi$ changes the Lagrangian density
only by a total derivative, $\Lag \to \Lag +
\epsilon\,\partial_\mu K^\mu$ (so the action is invariant for
suitable boundary conditions).  Then the \textbf{Noether current}
%
\begin{equation}
  j^\mu = \frac{\partial\Lag}{\partial(\partial_\mu\phi)}\,\Delta\phi
          - K^\mu
  \label{eq:noether_current}
\end{equation}
%
is conserved on solutions of the equations of motion:
%
\begin{equation}
  \partial_\mu j^\mu = 0
  \qquad\Longrightarrow\qquad
  Q = \int d^3\bx\, j^0 \;\;\text{is constant in time.}
\end{equation}
%
Proof sketch: compute $\delta\Lag$ under the transformation using the
chain rule, substitute the Euler--Lagrange equation
\eqref{eq:field_EL} to eliminate $\partial\Lag/\partial\phi$, and the
result collapses to $\partial_\mu(\text{current}) =
\partial_\mu K^\mu$.

Continuous symmetry $\Rightarrow$ conservation law.  This is the
deepest structural fact in field theory: energy, momentum, angular
momentum, charge, and (for FDM) particle number are all Noether
charges.

\subsection{Spacetime translations: the stress-energy tensor}
\label{sec:stress}

Invariance under translations $x^\nu \to x^\nu + \epsilon^\nu$
(one symmetry per direction $\nu$, with $\Delta\phi = \partial_\nu\phi$
and $K^\mu = \delta^\mu_\nu\,\Lag$) yields four conserved currents,
assembled into the \textbf{canonical stress-energy tensor}
%
\begin{equation}
  T^\mu{}_\nu
    = \frac{\partial\Lag}{\partial(\partial_\mu\phi)}\,\partial_\nu\phi
      - \delta^\mu_\nu\,\Lag,
  \qquad
  \partial_\mu T^\mu{}_\nu = 0.
  \label{eq:stress_tensor}
\end{equation}
%
The $\nu=0$ charge is the energy; its density for the free scalar is
%
\begin{equation}
  T^0{}_0
    = \dot\phi\,\dot\phi - \Lag
    = \tfrac{1}{2}\dot\phi^2 + \tfrac{1}{2}(\nabla\phi)^2
      + \tfrac{1}{2}m^2\phi^2
    \;\equiv\; \rho_\phi,
  \label{eq:energy_density}
\end{equation}
%
and the isotropic pressure (from the spatial diagonal, for a
homogeneous field) is
%
\begin{equation}
  p_\phi = \tfrac{1}{2}\dot\phi^2 - \tfrac{1}{2}m^2\phi^2
  \qquad (\nabla\phi = 0).
  \label{eq:pressure}
\end{equation}
%
These two lines are the engine of the misalignment story in
\texttt{fdm.tex} \S2: a frozen field ($\dot\phi\simeq 0$) has
$p_\phi = -\rho_\phi$, i.e.\ $w=-1$, while a field oscillating at
frequency $m$ has $\langle\tfrac{1}{2}\dot\phi^2\rangle =
\langle\tfrac{1}{2}m^2\phi^2\rangle$ (virial theorem for a harmonic
oscillator), so $\langle p_\phi\rangle = 0$ and the field redshifts
like pressureless matter.

(A convention flag: in the GR literature the stress tensor is defined
by varying the action with respect to $g^{\mu\nu}$ --- the ``Hilbert''
stress tensor --- and for a perfect fluid
$T^\mu{}_\nu = \mathrm{diag}(-\rho,p,p,p)$, so $T^0{}_0 = -\rho$
there.  The canonical Noether tensor above differs by an overall
sign convention in the $0$-components; the physics --- which
combination of $\dot\phi^2$, $(\nabla\phi)^2$, $m^2\phi^2$ is the
energy density --- is the same.  See \S\ref{sec:conventions}.)

\subsection{Internal symmetry: the complex scalar and particle number}
\label{sec:u1}

A \emph{complex} scalar field with action
%
\begin{equation}
  S = \int d^4x \left[
      -g^{\mu\nu}\partial_\mu\phi^*\,\partial_\nu\phi
      - m^2|\phi|^2 \right]
  \label{eq:complex_action}
\end{equation}
%
is invariant under the global $U(1)$ phase rotation
$\phi \to e^{-i\alpha}\phi$.  The Noether charge is
%
\begin{equation}
  Q = i\!\int d^3\bx\,
      \bigl(\phi^*\dot\phi - \phi\,\dot\phi^*\bigr),
  \label{eq:u1_charge}
\end{equation}
%
(sign chosen so that $Q>0$ for particles).  In the non-relativistic
limit, write $\phi = e^{-imt}\,\psi/\sqrt{2m}$ with $\psi$ slowly
varying; then $\dot\phi \simeq -im\phi$ and
%
\begin{equation}
  Q \;\longrightarrow\; \int d^3\bx\, |\psi|^2 \;=\; N,
\end{equation}
%
the \textbf{particle number}.  The conserved $U(1)$ charge of the
relativistic theory becomes the conservation of particle number in
the Schr\"odinger description, and $\rho = m|\psi|^2$ is the mass
density.

A \emph{real} scalar (the axion) has no exact $U(1)$ symmetry and no
exactly conserved number.  But in the non-relativistic regime,
number-changing processes require energies of order $m$ per quantum
and are suppressed by powers of the (tiny) couplings and of
$(v/c)^2$; particle number is conserved to superb accuracy.  That is
why \texttt{fdm.tex} can package the real field's slow envelope into
a complex Schr\"odinger field $\psi$ and treat $\int|\psi|^2$ as
constant.

%----------------------------------------------------------------------
\section{Hamiltonian formulation}
\label{sec:hamiltonian}

Quantisation needs the Hamiltonian.  Define the \textbf{canonical
momentum density} conjugate to $\phi$,
%
\begin{equation}
  \pi(\bx,t) \equiv \frac{\partial\Lag}{\partial\dot\phi} = \dot\phi
  \qquad\text{(free scalar)},
\end{equation}
%
and Legendre-transform:
%
\begin{equation}
  \Ham = \pi\dot\phi - \Lag
       = \tfrac{1}{2}\pi^2 + \tfrac{1}{2}(\nabla\phi)^2
         + \tfrac{1}{2}m^2\phi^2,
  \qquad
  H = \int d^3\bx\,\Ham.
  \label{eq:hamiltonian}
\end{equation}
%
This is Eq.~\eqref{eq:energy_density} again --- the Hamiltonian is
the Noether energy.  Hamilton's equations
$\dot\phi = \delta H/\delta\pi$, $\dot\pi = -\delta H/\delta\phi$
reproduce the KG equation.

In a periodic box of volume $V$ (the natural setting for both lattice
simulations and quantisation bookkeeping), expand in Fourier modes
$\phi(\bx) = V^{-1/2}\sum_\bk \phi_\bk\, e^{i\bk\cdot\bx}$ with
$\phi_{-\bk} = \phi_\bk^*$ (reality).  The Hamiltonian separates:
%
\begin{equation}
  H = \sum_{\bk} \left[
        \tfrac{1}{2}|\pi_\bk|^2
        + \tfrac{1}{2}\omega_k^2\,|\phi_\bk|^2
      \right],
  \qquad
  \omega_k^2 = k^2 + m^2.
  \label{eq:H_modes}
\end{equation}
%
\textbf{The free field is exactly a collection of independent harmonic
oscillators, one per Fourier mode $\bk$, with frequency $\omega_k$.}
This single sentence is the entire content of free-field quantisation;
the next section merely applies the known quantum mechanics of the
harmonic oscillator to each mode.

%----------------------------------------------------------------------
\section{Canonical quantisation}
\label{sec:quantization}

\subsection{One oscillator}
\label{sec:oscillator}

Recall the quantum harmonic oscillator with
$H = \tfrac{1}{2}p^2 + \tfrac{1}{2}\omega^2 q^2$ and $[q,p]=i$.
Define ladder operators
%
\begin{equation}
  a = \sqrt{\tfrac{\omega}{2}}\left(q + \tfrac{i}{\omega}p\right),
  \qquad
  a^\dagger = \sqrt{\tfrac{\omega}{2}}\left(q - \tfrac{i}{\omega}p\right),
  \qquad
  [a, a^\dagger] = 1,
\end{equation}
%
in terms of which $H = \omega\,(a^\dagger a + \tfrac{1}{2})$.  The
spectrum is $E_n = \omega\,(n + \tfrac{1}{2})$, $n = 0,1,2,\dots$;
$a^\dagger$ raises $n$ by one, $a$ lowers it, and
$\hat n = a^\dagger a$ counts quanta.

\subsection{Quantising every mode}
\label{sec:field_quant}

Promote $\phi(\bx)$ and $\pi(\bx)$ to operators with the equal-time
commutator
%
\begin{equation}
  [\hat\phi(\bx), \hat\pi(\bx')] = i\,\delta^3(\bx - \bx'),
\end{equation}
%
the continuum version of $[q,p]=i$ at every point.  Because the
Hamiltonian \eqref{eq:H_modes} is a sum of independent oscillators,
introduce one pair $(a_\bk, a^\dagger_\bk)$ per mode with
$[a_\bk, a^\dagger_{\bk'}] = \delta_{\bk\bk'}$:
%
\begin{equation}
  \hat\phi(\bx)
    = \frac{1}{\sqrt V}\sum_\bk \frac{1}{\sqrt{2\omega_k}}
      \left( a_\bk\, e^{i\bk\cdot\bx}
           + a^\dagger_\bk\, e^{-i\bk\cdot\bx} \right),
  \qquad
  H = \sum_\bk \omega_k \left( a^\dagger_\bk a_\bk
        + \tfrac{1}{2} \right).
  \label{eq:mode_expansion}
\end{equation}
%
(The infinite-volume convention replaces $V^{-1/2}\sum_\bk \to
\int\! d^3k/(2\pi)^3$ and $\delta_{\bk\bk'} \to
(2\pi)^3\delta^3(\bk-\bk')$; see \S\ref{sec:conventions}.)

The structure of the theory is now read off:
%
\begin{itemize}
  \item \textbf{Vacuum} $|0\rangle$: all modes in their ground state,
        $a_\bk|0\rangle = 0$.
  \item \textbf{Particles}: $a^\dagger_\bk|0\rangle$ is a state of
        energy $\omega_k = \sqrt{k^2+m^2}$ and momentum $\bk$ --- one
        quantum of the field, i.e.\ one particle of mass $m$.
        ``Particles'' are excitations of the field, not ingredients
        added to it.
  \item \textbf{Fock space}: repeated $a^\dagger$'s build
        $|n_{\bk_1}, n_{\bk_2}, \dots\rangle$ with any number of
        quanta per mode.  Since the $a^\dagger$'s commute, the states
        are automatically symmetric: a scalar field describes
        \textbf{bosons}, with no limit on the occupation
        $n_\bk = \langle a^\dagger_\bk a_\bk\rangle$ of a single
        mode.  This is what makes the condensed, classical-wave
        regime of \S\ref{sec:classical} possible.
  \item \textbf{Vacuum energy}: the $\sum_\bk \tfrac{1}{2}\omega_k$
        zero-point term diverges.  In non-gravitational physics only
        energy differences matter and the constant is discarded
        (normal ordering).  Gravity does couple to it --- that
        unresolved mismatch is the cosmological-constant problem,
        outside our scope.
\end{itemize}

For the complex scalar, the same construction produces two sets of
operators $(a_\bk, b_\bk)$ --- particles and antiparticles --- and
the $U(1)$ charge \eqref{eq:u1_charge} becomes
$Q = \sum_\bk (a^\dagger_\bk a_\bk - b^\dagger_\bk b_\bk)$, the
particle-minus-antiparticle count.

\subsection{What we skipped}
\label{sec:skipped}

A full QFT course continues with propagators, interacting-field
perturbation theory, Feynman diagrams, and renormalisation
\citep{PeskinSchroeder1995, Srednicki2007}.  None of that machinery
is needed for FDM cosmology: the field is free to excellent
approximation (\S\ref{sec:interactions}), and the questions asked ---
how does the classical wave gravitate and interfere? --- live
entirely in the sector described above.

%----------------------------------------------------------------------
\section{Coherent states and the classical wave limit}
\label{sec:classical}

\subsection{When is a quantum field classical?}
\label{sec:when_classical}

The classical field equation (KG, or its Schr\"odinger reduction) is
exact for the operator $\hat\phi$ in the Heisenberg picture of the
free theory.  The physical question is whether the \emph{state} is
such that $\hat\phi$ behaves like a definite classical configuration:
expectation value large compared to quantum fluctuations.

The answer is controlled by the occupation number per mode.  A mode
with $n$ quanta has field amplitude $\propto\sqrt{n}$ while its
quantum fluctuation stays at the zero-point level, so the fractional
``quantumness'' scales as $1/\sqrt{n}$.

\subsection{Coherent states}
\label{sec:coherent}

The precise statement uses \textbf{coherent states}, the minimum-
uncertainty states of the harmonic oscillator:
%
\begin{equation}
  |\alpha\rangle = e^{-|\alpha|^2/2}\, e^{\alpha a^\dagger}\,
                   |0\rangle,
  \qquad
  a\,|\alpha\rangle = \alpha\,|\alpha\rangle,
  \qquad \alpha \in \mathbb{C}.
\end{equation}
%
Key properties:
%
\begin{itemize}
  \item $\langle \hat n\rangle = |\alpha|^2$, with Poisson number
        statistics: $\Delta n/\langle n\rangle = 1/|\alpha| \to 0$
        for large occupation.
  \item $\langle\alpha|\hat q(t)|\alpha\rangle$ oscillates exactly
        like the classical trajectory with amplitude
        $\propto|\alpha|$; the quantum spread around it is the fixed
        zero-point width, fractionally negligible for
        $|\alpha|\gg 1$.
  \item A classical driving force acting on the vacuum produces
        precisely a coherent state.  In the field theory, a slowly
        varying classical source --- e.g.\ the misalignment initial
        condition, a homogeneous field displacement $\phi_i$ across
        the horizon --- populates the $\bk\simeq 0$ mode into a
        coherent state of enormous $|\alpha|$.
\end{itemize}
%
A quantum field in a multi-mode coherent state
$\prod_\bk|\alpha_\bk\rangle$ has
$\langle\hat\phi(\bx,t)\rangle = \phi_{\rm cl}(\bx,t)$, where
$\phi_{\rm cl}$ solves the \emph{classical} field equation with mode
amplitudes $\alpha_\bk$.  Relative corrections are
$O(1/\sqrt{n_\bk})$.

\subsection{FDM occupation numbers}
\label{sec:occupation}

For dark matter of density $\rho$ made of quanta of mass $m$ moving
at velocity $v$, the occupation of the relevant modes is the number
of quanta per de Broglie volume:
%
\begin{equation}
  n_{\rm occ} \sim \frac{\rho}{m}\,
      \left(\frac{2\pi\hbar}{m v}\right)^{3}.
\end{equation}
%
With $\rho \sim 0.3\;\mathrm{GeV\,cm^{-3}}$ (local DM density),
$m = 10^{-22}\eV$, $v \sim 100\kms$:
%
\begin{equation}
  \frac{\rho}{m} \sim 3\times10^{30}\,\mathrm{cm^{-3}},
  \qquad
  \lambda_{\rm dB}^3 \sim 5\times10^{64}\,\mathrm{cm^{3}}
  \qquad\Longrightarrow\qquad
  n_{\rm occ} \sim 10^{95}.
\end{equation}
%
Compare a WIMP with $m = 100\,$GeV: $n_{\rm occ} \sim 10^{-28}$ ---
occupation so far below unity that the particle description is the
only sensible one.  FDM sits at the opposite extreme: $10^{95}$
quanta per mode make the coherent-state description essentially
exact, with quantum corrections of order $10^{-47}$.  \textbf{The
``classical wave'' treatment of FDM is not an approximation scheme
that might break down --- it is more accurate than any other
approximation in cosmology.}  The wave-like phenomenology
(interference granules, solitons, quantum pressure) is classical
wave mechanics; ``quantum'' enters only through $\hbar/m$ setting
the wavelength \citep{Hui2017}.

Caveat for precision: over very long times, mode couplings from
gravitational interactions can degrade coherence (``quantum
breaking''); estimates place the timescale far beyond the Hubble time
for FDM parameters, so cosmological simulations solving the
classical Schr\"odinger--Poisson system are safe.

%----------------------------------------------------------------------
\section{Interactions and the axion potential}
\label{sec:interactions}

\subsection{Adding a potential}
\label{sec:potential}

Interactions replace $\tfrac{1}{2}m^2\phi^2 \to V(\phi)$.  The
Euler--Lagrange equation becomes $\Box\phi = V'(\phi)$.  The simplest
interacting theory adds a quartic term,
%
\begin{equation}
  V(\phi) = \tfrac{1}{2}m^2\phi^2 + \tfrac{\lambda}{4!}\phi^4,
\end{equation}
%
which in the quantum theory makes quanta scatter off each other with
amplitude $\propto\lambda$.

\subsection{The axion cosine}
\label{sec:axion}

Axion-like particles are born as angular fields: $\phi = f\theta$
with $\theta$ a periodic phase and $f$ the \textbf{decay constant}.
Non-perturbative effects (instantons) generate the periodic
potential
%
\begin{equation}
  V(\phi) = m^2 f^2 \left[ 1 - \cos\!\left(\frac{\phi}{f}\right)
            \right]
          = \tfrac{1}{2}m^2\phi^2
            - \frac{m^2}{4!\,f^2}\,\phi^4 + \dots
  \label{eq:axion_potential}
\end{equation}
%
Two structural consequences:
%
\begin{itemize}
  \item \textbf{Why the mass is naturally tiny.}  Before the
        instanton effects switch on, the action is invariant under
        the \emph{shift symmetry} $\phi \to \phi + \mathrm{const}$
        (only $\partial_\mu\phi$ appears), which forbids any
        potential at all.  The mass is generated entirely by
        symmetry-breaking effects of size $m \sim M_{\rm Pl}\,
        e^{-S_{\rm inst}}$, exponentially suppressed by an instanton
        action \citep{Arvanitaki2010}.  A tiny $m$ is therefore
        \emph{technically natural}: quantum corrections cannot drag
        it upward, because every correction must itself break the
        shift symmetry and carries the same exponential suppression.
        (Contrast the Higgs, whose light mass is radiatively
        unstable --- the hierarchy problem.)  This is why theorists
        take $m \sim 10^{-22}\eV$ seriously rather than dismissing
        it as grotesque fine-tuning.
  \item \textbf{Why self-interactions are negligible.}  Expanding the
        cosine gives an \emph{attractive} quartic coupling
        $\lambda = -m^2/f^2$.  For FDM parameters
        ($m = 10^{-22}\eV$, $f \sim 10^{17}\,$GeV $= 10^{26}\eV$):
        %
        \begin{equation}
          |\lambda| = \frac{m^2}{f^2}
            = \frac{10^{-44}}{10^{52}} = 10^{-96}.
        \end{equation}
        %
        The most weakly self-coupled matter ever contemplated.
        Self-interactions matter only where the field amplitude
        approaches $f$ (e.g.\ possibly in the densest soliton
        cores); everywhere else the free-field treatment of
        \S\ref{sec:free}--\ref{sec:classical} applies
        \citep{Hui2017, Marsh2016}.
\end{itemize}

%----------------------------------------------------------------------
\section{The scalar field on curved spacetime}
\label{sec:curved}

\subsection{Minimal coupling}
\label{sec:minimal}

To put the field on a curved background, apply the ``comma-goes-to-
semicolon'' rule: replace $\eta_{\mu\nu} \to g_{\mu\nu}(x)$ and
$d^4x \to d^4x\sqrt{-g}$ in the action --- which is exactly
Eq.~\eqref{eq:scalar_action}.  This is \emph{minimal coupling}: the
field senses geometry only through the metric appearing in its
kinetic term and volume element.

Varying the action now gives the covariant KG equation
%
\begin{equation}
  \Box_g\phi \equiv \frac{1}{\sqrt{-g}}\,
    \partial_\mu\!\bigl(\sqrt{-g}\,g^{\mu\nu}\partial_\nu\phi\bigr)
    = V'(\phi).
  \label{eq:KG_covariant}
\end{equation}
%
On flat FRW ($\sqrt{-g} = a^3$, $g^{00}=-1$, $g^{ii}=a^{-2}$) with
$V = \tfrac{1}{2}m^2\phi^2$:
%
\begin{equation}
  \ddot\phi + 3H\dot\phi - \frac{1}{a^2}\nabla^2\phi + m^2\phi = 0.
  \label{eq:KG_FRW}
\end{equation}
%
The $3H\dot\phi$ Hubble-friction term is bookkeeping from
$\sqrt{-g}=a^3$: the physical energy in a comoving volume dilutes as
the volume grows.  Dropping gradients (homogeneous background field)
gives the misalignment equation of \texttt{fdm.tex} \S2; keeping them
and taking the non-relativistic limit gives Schr\"odinger--Poisson
(\texttt{fdm.tex} \S3).

\subsection{Non-minimal coupling}
\label{sec:nonminimal}

Curvature admits one more renormalisable term:
%
\begin{equation}
  \Delta\Lag = -\tfrac{1}{2}\,\xi\,R\,\phi^2,
\end{equation}
%
with $R$ the Ricci scalar.  $\xi = 0$ is minimal coupling;
$\xi = 1/6$ is ``conformal'' coupling (the massless theory then
ignores conformally flat curvature).  Non-minimally coupled scalars
drive Higgs inflation and related models, but the FDM literature
uniformly assumes $\xi = 0$: the field just rides the geometry.

Quantum field theory \emph{on} curved spacetime (particle creation,
Hawking/Unruh effects, the ambiguity of the vacuum when the
background is time-dependent) is a rich subject
\citep{BirrellDavies1982}; for FDM at $z \lesssim 10^6$ the modes of
interest are so far inside the horizon and so highly occupied that
none of these subtleties surface.

%----------------------------------------------------------------------
\section{Convention traps}
\label{sec:conventions}

Formulas imported from different references disagree by signs and
factors.  The four recurring traps:
%
\begin{itemize}
  \item \textbf{Metric signature.}  ``Mostly plus'' $(-,+,+,+)$:
        most GR and cosmology texts (Wald, Weinberg, MTW, and these
        notes).  ``Mostly minus'' $(+,-,-,-)$: most particle-physics
        texts \citep{PeskinSchroeder1995}.  The kinetic term flips
        written sign between the two,
        $-\tfrac{1}{2}g^{\mu\nu}\partial_\mu\phi\partial_\nu\phi
        \leftrightarrow
        +\tfrac{1}{2}g^{\mu\nu}\partial_\mu\phi\partial_\nu\phi$,
        while describing the same physics.
  \item \textbf{Factors of $\tfrac{1}{2}$ and canonical
        normalisation.}  The $\tfrac{1}{2}$ prefactors give the
        field a canonically normalised kinetic term; complex fields
        drop them (Eq.~\eqref{eq:complex_action}) because
        $\phi$ and $\phi^*$ are independent.  Rescalings
        $\phi \to c\phi$ move constants into masses and couplings.
  \item \textbf{Fourier and mode-normalisation conventions.}
        Box normalisation ($V^{-1/2}\sum_\bk$,
        $[a_\bk,a^\dagger_{\bk'}]=\delta_{\bk\bk'}$) versus
        continuum ($\int d^3k/(2\pi)^3$,
        $[a_\bk,a^\dagger_{\bk'}]=(2\pi)^3\delta^3(\bk-\bk')$),
        and relativistic normalisation factors $\sqrt{2\omega_k}$
        placed differently by different books.  Physical quantities
        (energies, occupation numbers) agree.
  \item \textbf{Stress-tensor sign.}  The canonical Noether tensor
        \eqref{eq:stress_tensor} has $T^0{}_0 = +\rho$; the Hilbert
        tensor used in GR has $T^0{}_0 = -\rho$ in mostly-plus
        signature.  Check which one a reference means before
        matching components.
\end{itemize}

%----------------------------------------------------------------------
\section{Summary}
\label{sec:summary}

The chain of ideas, from mechanics to FDM cosmology:

\begin{center}
\begin{tabular}{@{}lll@{}}
\toprule
Mechanics & Classical field theory & Quantum field theory \\
\midrule
Coordinate $q(t)$ & Field $\phi(\bx,t)$ & Operator $\hat\phi(\bx,t)$ \\
$L(q,\dot q)$ & Density $\Lag(\phi,\partial_\mu\phi)$ & same $\Lag$, quantised \\
$S = \int L\,dt$ & $S = \int d^4x\sqrt{-g}\,\Lag$ & path-integral weight \\
Newton from $\delta S{=}0$ & Klein--Gordon from $\delta S{=}0$ & Heisenberg equations \\
Energy $E$ conserved & Noether currents $T^\mu{}_\nu$, $j^\mu$ & $\hat H$, $\hat Q = \hat N$ \\
Oscillator modes & Fourier modes $\omega_k^2 = k^2{+}m^2$ & quanta = particles \\
Amplitude & Mode amplitude $\phi_\bk$ & occupation $n_\bk$ \\
--- & classical wave & coherent state, $n_\bk \gg 1$ \\
\bottomrule
\end{tabular}
\end{center}

For FDM the essential facts are: (i) the action
\eqref{eq:scalar_action} is the generic low-energy description of any
light scalar, with the tiny mass protected by a shift symmetry;
(ii) its quanta are bosons, so occupation numbers are unbounded;
(iii) misalignment populates the low-$k$ modes to
$n_{\rm occ} \sim 10^{95}$, making the coherent-state / classical-wave
description essentially exact; and (iv) everything in
\texttt{fdm.tex} --- Schr\"odinger--Poisson, quantum pressure,
solitons, interference granules --- is then the \emph{classical}
dynamics of that wave.

\section*{Further reading}

\citet{Tong2007} is the gentlest complete introduction to canonical
quantisation; \citet{Coleman2011} is the classic lecture course with
the most physical insight per page; \citet{Srednicki2007} is a
systematic textbook whose draft is freely available;
\citet{PeskinSchroeder1995} is the standard reference for
perturbative calculations.  For the FDM application,
\citet{Hui2017} \S2 covers the ground of these notes at research
level, and \citet{Marsh2016} reviews axion cosmology broadly.
Local copies of the freely available texts are kept in
\texttt{\textasciitilde/Documents/ref\_qft/}.

\bibliography{references}

\end{document}
